% ============================================================
% CHAPITRE 9 — Industrialisation DevOps
% ============================================================
\chapter{Industrialisation DevOps --- Docker, Jenkins et SonarQube}

\section{Introduction}

Après la réalisation fonctionnelle de la plateforme (microservices
\texttt{EUREKA}, \texttt{GATEWAY}, \texttt{USER}, \texttt{RECRUITMENT},
\texttt{RECLAMATION} et frontend Angular), ce chapitre décrit
l'\textbf{industrialisation} de la chaîne de livraison. L'objectif est de
passer d'une exécution manuelle locale à un cycle automatisé
\textbf{CI/CD}~: construction, analyse qualité, publication d'images Docker
et déploiement des conteneurs --- \textbf{backend et frontend inclus}.

La chaîne mise en place s'appuie sur trois briques complémentaires~:
\begin{itemize}
  \item \textbf{Docker} pour empaqueter chaque microservice et le frontend
        de façon reproductible~;
  \item \textbf{Jenkins} pour orchestrer les pipelines CI (intégration
        continue) et CD (déploiement)~;
  \item \textbf{SonarQube} pour contrôler la qualité du code avant
        publication de l'image.
\end{itemize}

Le code source est hébergé sur \textbf{GitHub}~; Jenkins clone le dépôt,
exécute les étapes définies dans des \texttt{Jenkinsfile} versionnés sous
\texttt{devops/pipelines/}, puis pousse les images vers
\textbf{Docker Hub} (compte \texttt{lfray}).


\section{Objectifs}

\begin{itemize}
  \item conteneuriser chaque microservice backend avec un
        \texttt{Dockerfile} multi-étapes (Maven~17 puis JRE~17)~;
  \item conteneuriser le \textbf{frontend Angular} (build Node puis
        serveur \textbf{nginx})~;
  \item automatiser, pour chaque composant~: \emph{checkout},
        build (\emph{Maven} ou \emph{npm}), analyse \emph{SonarQube},
        \emph{Quality Gate}, \emph{docker build} et \emph{docker push}~;
  \item déployer chaque composant via un job CD~: \emph{docker pull} puis
        \emph{docker run}~;
  \item garantir la cohérence des environnements (profil Spring
        \texttt{docker}, \texttt{environment.docker.ts} pour Angular).
\end{itemize}


\section{Conteneurisation Docker}

\subsection{Principe du Dockerfile multi-étapes}

Chaque microservice dispose d'un \texttt{Dockerfile} calqué sur le même
modèle~:
\begin{enumerate}
  \item étape \textbf{build}~: image \texttt{maven:3.9.9-eclipse-temurin-17},
        compilation et packaging du JAR (\texttt{mvn package -DskipTests})~;
  \item étape \textbf{runtime}~: image légère
        \texttt{eclipse-temurin:17-jre-jammy}, copie du JAR, utilisateur non
        privilégié, exposition du port métier.
\end{enumerate}

Cette approche réduit la taille de l'image finale et sépare clairement les
outils de compilation de l'environnement d'exécution.

\subsection{Frontend Angular (Node + nginx)}

Le frontend suit un schéma multi-étapes adapté au web~:
\begin{enumerate}
  \item étape \textbf{build}~: image \texttt{node:20-alpine},
        \texttt{npm ci} puis \texttt{ng build --configuration=docker}
        (fichier \texttt{environment.docker.ts} pointant l'API vers la
        Gateway sur \texttt{http://localhost:8088})~;
  \item étape \textbf{runtime}~: image \texttt{nginx:1.27-alpine} servant
        les fichiers statiques du dossier \texttt{dist/frontend/browser},
        avec une configuration SPA (\texttt{try\_files} vers
        \texttt{index.html}).
\end{enumerate}

Les jobs Jenkins associés sont \texttt{ci-frontend} et
\texttt{cd-frontend}~; l'interface est exposée sur
\texttt{http://localhost:4200}.

\subsection{Profil Spring \texttt{docker}}

Un fichier \texttt{application-docker.properties} adapte la configuration
lorsque \texttt{SPRING\_PROFILES\_ACTIVE=docker}~:
\begin{itemize}
  \item URL Eureka~: \texttt{http://host.docker.internal:8761/eureka}
        (accès au registre depuis un conteneur sur Docker Desktop)~;
  \item URL JDBC MySQL pointant vers \texttt{host.docker.internal:3306}
        pour les services \texttt{USER}, \texttt{RECRUITMENT} et
        \texttt{RECLAMATION}.
\end{itemize}

\subsection{Images produites}

\begin{table}[htbp]
  \centering
  \small
  \caption{Images Docker Hub de la plateforme DAAM}
  \label{tab:docker-images}
  \begin{tabular}{|l|l|c|}
    \hline
    \textbf{Service} & \textbf{Image} & \textbf{Port conteneur} \\
    \hline
    Eureka & \texttt{lfray/daam-eureka} & 8761 \\
    Gateway & \texttt{lfray/daam-gateway} & 8080 \\
    User & \texttt{lfray/daam-user} & 8089 \\
    Recruitment & \texttt{lfray/daam-recruitment} & 8090 \\
    Reclamation & \texttt{lfray/daam-reclamation} & 8091 \\
    Frontend & \texttt{lfray/daam-frontend} & 80 \\
    \hline
  \end{tabular}
\end{table}

La Gateway est exposée sur le port hôte \textbf{8088} (mapping
\texttt{8088:8080}) afin d'éviter le conflit avec Jenkins, déjà écouté sur
le port~8080. Le frontend est exposé sur le port hôte \textbf{4200}
(mapping \texttt{4200:80}).

\begin{figure}[H]
  \centering
  \reportfig{img/devops-docker-containers.png}
  \caption{Docker Desktop --- conteneurs de la plateforme DAAM
           (Jenkins, SonarQube et microservices déployés)}
  \label{fig:devops-docker-containers}
\end{figure}


\section{Intégration continue avec Jenkins}

\subsection{Organisation des pipelines}

Pour chaque microservice, deux jobs Pipeline sont définis~:
\begin{itemize}
  \item \textbf{CI} (\texttt{ci-<service>})~: qualité + construction +
        publication~;
  \item \textbf{CD} (\texttt{cd-<service>})~: récupération de l'image et
        démarrage du conteneur.
\end{itemize}

Les scripts sont versionnés dans le dépôt, par exemple~:
\texttt{devops/pipelines/ci-eureka.Jenkinsfile} et
\texttt{cd-eureka.Jenkinsfile}. Jenkins est configuré en mode
\emph{Pipeline script from SCM} (branche \texttt{master}, clone peu
profond).

\subsection{Étapes du pipeline CI}

Le pipeline CI suit le déroulé suivant~:
\begin{enumerate}
  \item \textbf{Checkout} GitHub (authentification par jeton)~;
  \item \textbf{Build}~: \emph{Maven clean package} (backend) ou
        \emph{npm ci / ng build} (frontend), exécuté dans un conteneur~;
  \item \textbf{SonarQube} via \texttt{sonar-scanner-cli}~;
  \item \textbf{Quality Gate} (\texttt{waitForQualityGate})~;
  \item \textbf{Docker build} du composant~;
  \item \textbf{Docker push} vers Docker Hub
        (\texttt{lfray/daam-<composant>:latest}).
\end{enumerate}

Maven et Sonar s'exécutent eux-mêmes dans Docker afin d'éviter
l'installation de JDK/Maven dans l'agent Jenkins. Le volume nommé
\texttt{jenkins\_home} est monté explicitement, car un simple
\texttt{-v \$PWD:/app} est incorrect lorsque Jenkins tourne déjà dans un
conteneur (problème classique \emph{Docker-from-Docker} sous Docker
Desktop).

\subsection{Étapes du pipeline CD}

Le pipeline CD est volontairement simple~:
\begin{enumerate}
  \item connexion Docker Hub~;
  \item \texttt{docker pull} de l'image publiée par le CI~;
  \item \texttt{docker rm -f} de l'ancien conteneur (s'il existe)~;
  \item \texttt{docker run -d} avec le profil \texttt{docker} et le
        mapping de ports~;
  \item déconnexion Docker Hub.
\end{enumerate}

\subsection{Ordre de déploiement recommandé}

\begin{enumerate}
  \item CI puis CD d'\textbf{Eureka} (registre)~;
  \item CI puis CD de la \textbf{Gateway} et des microservices métier
        (\texttt{USER}, \texttt{RECRUITMENT}, \texttt{RECLAMATION}), une
        fois Eureka disponible~;
  \item CI puis CD du \textbf{frontend} (après Gateway, pour que l'API
        soit joignable)~;
  \item MySQL doit être accessible sur l'hôte (\texttt{localhost:3306})
        pour les services persistants.
\end{enumerate}

\begin{figure}[H]
  \centering
  \reportfig{img/devops-jenkins-pipelines.png}
  \caption{Jenkins --- vue d'ensemble des pipelines CI/CD des microservices DAAM}
  \label{fig:devops-jenkins-pipelines}
\end{figure}


\section{Contrôle qualité avec SonarQube}

Chaque analyse CI crée (ou met à jour) un projet Sonar dédié
(\texttt{daam-eureka}, \texttt{daam-gateway}, \texttt{daam-frontend},
etc.). Le scanner analyse les sources Java (backend) ou TypeScript /
HTML (frontend). Le Quality Gate \emph{Sonar way} doit être
\textbf{Passed} avant d'autoriser le build et le push Docker.

Jenkins communique avec SonarQube via l'URL
\texttt{http://host.docker.internal:9000}. Un \textbf{webhook} Sonar
vers Jenkins
(\texttt{http://host.docker.internal:8080/sonarqube-webhook/})
permet au stage \texttt{waitForQualityGate} de se débloquer dès que
l'analyse serveur est terminée.

\begin{figure}[H]
  \centering
  \reportfig{img/devops-sonar-projects.png}
  \caption{SonarQube --- projets DAAM et résultats du Quality Gate}
  \label{fig:devops-sonar-projects}
\end{figure}


\section{Environnement technique retenu}

\begin{table}[htbp]
  \centering
  \small
  \caption{Stack DevOps utilisée pour DAAM}
  \label{tab:devops-stack}
  \begin{tabular}{|l|p{8.5cm}|}
    \hline
    \textbf{Composant} & \textbf{Rôle} \\
    \hline
    GitHub & Gestion de versions et SCM des pipelines \\
    Docker Desktop & Moteur de conteneurs (hôte Windows) \\
    Jenkins (conteneur) & Orchestration CI/CD \\
    SonarQube (conteneur) & Analyse statique et Quality Gate \\
    Docker Hub & Registre d'images (\texttt{lfray}) \\
    Maven 3.9 / JDK 17 & Compilation des microservices Java \\
    Node.js 20 / Angular 18 & Build du frontend \\
    nginx & Serveur web des fichiers Angular (image frontend) \\
    \hline
  \end{tabular}
\end{table}


\section{Apports et limites}

\subsection{Apports}
\begin{itemize}
  \item reproductibilité des builds (même Dockerfile, mêmes images)~;
  \item détection précoce des défauts via SonarQube~;
  \item déploiement homogène service par service~;
  \item séparation claire CI (qualité + artefact) / CD (exécution).
\end{itemize}

\subsection{Limites actuelles}
\begin{itemize}
  \item pas encore de \texttt{docker-compose} unifié pour démarrer
        toute la stack d'un seul ordre (déploiement job par job)~;
  \item MySQL reste installé sur l'hôte (hors conteneur CI/CD)~;
  \item couverture de tests encore faible (visible dans SonarQube).
\end{itemize}


\section{Conclusion}

Ce chapitre a présenté la mise en place d'une chaîne DevOps adaptée à
l'architecture de DAAM~: conteneurisation Docker, pipelines Jenkins CI/CD
par composant (microservices \textbf{et} frontend Angular) et contrôle
qualité SonarQube. L'ensemble peut ainsi être construit, analysé, publié
puis déployé de façon reproductible.

La \textbf{conclusion générale} du rapport synthétise l'ensemble du
projet (besoins, conception, sprints et industrialisation) et ouvre
quelques perspectives d'évolution.

\clearpage
